Let $\Delta Q_{ikt}$ be candidate-attributable incremental quantity in benefit
category $k$ for OD relation $i$ and appraisal year $t$, and let $v_{kt}$ be the
corresponding real unit value in the common price base. Annual benefit is
\begin{equation}
B_t=\sum_{i\in\ODset}\sum_k \Delta Q_{ikt}v_{kt}.
\label{eq:annual-benefit}
\end{equation}
For bicycle type $b\in\{\mathrm{conventional},\mathrm{e\mbox{-}bike}\}$, let
$\Delta K_{bt}$ be eligible incremental cycle-kilometres, $U_{bt}$ incremental
users, $v_b^H$ the health value per kilometre, and $H_b^{\max}$ the annual cap
per user. Type-specific health benefit is
\begin{equation}
B_{bt}^{H}
=\min\!\left(\Delta K_{bt}v_b^H,\ U_{bt}H_b^{\max}\right).
\label{eq:health-cap}
\end{equation}
Conventional and e-bike values and caps are applied separately before
aggregation and discounting; the same user-kilometre cannot enter both types.
For year-specific real discount rate $r_t$, define the cumulative discount
factor $D_0=1$ and $D_t=\prod_{u=1}^{t}(1+r_u)$. For horizon $T$, capital cost
$K_t$, operations and maintenance $M_t$, renewal $R_t$, and terminal residual
value $S_T$,
\begin{align}
\PV(B)
 &=\sum_{t=0}^{T}\frac{B_t}{D_t},
\label{eq:pv-benefit}\\
\PV(C)
 &=\sum_{t=0}^{T}\frac{K_t+M_t+R_t}{D_t}
 -\frac{S_T}{D_T},
\label{eq:pv-cost}\\
\BCR^{\mathrm{life}}
 &=\frac{\PV(B)}{\PV(C)}.
\label{eq:lifecycle-bcr}
\end{align}
For a non-commercial public-sector activity under NZTA General Circular 25/01,
the principal schedule is
\begin{equation}
r_t=
\begin{cases}
0.02, & 1\le t\le30,\\
0.015, & 31\le t\le100,\\
0.01, & t\ge101,
\end{cases}
\label{eq:discount-schedule}
\end{equation}
with a constant 8% sensitivity. The default workbench horizon is 40 years, so
only the first two principal-rate bands enter that calculation.

Opening-year demand ramp, asset life, renewals, residual value, price base, and
the applicable appraisal manual version must be declared. Benefit categories
are included only when their incremental quantities and values are compatible
and non-overlapping.
